\chapter{{\fontsize{14pt}{21pt}\selectfont\MakeUppercase{Заголовок 1}}}
\label{cha:analysis}

\section{Заголовок 2}


% Информация об источниках хранится здесь - rpz.bib, ее проще всего сначала собрать через zotero и потом экспортировать в bib.
Пример цитирования \cite{law125fz1} или \cite{who2025stats} или \cite{_laravel_2023} или \cite{nestjs-docs1}.


Пример создания таблиц~\ref{tab:role}, в которую были занесены основные аспекты каждой реализации.
% Перед \ref рекомендуется ставить тильду (~) чтобы цифру не отрывало.
% Для генерации таблиц можно использовать https://www.tablesgenerator.com/

\begin{table}[h]
\caption{Сравнительная таблица ролевых моделей}
\label{tab:role}
\begin{tabular}{|p{2cm}|p{2.5cm}|p{2.3cm}|p{4.5cm}|p{3cm}|}
\hline
\centering Модель &  \centering Подходит для CRM & \centering Гибкость &  \centering Масштабируемость & \multicolumn{1}{c|}{Сложность} \\ \hline
RBAC   & Да                         & Средняя  & Высокая          & Низкая    \\ \hline
ABAC   & Да                         & Высокая  & Высокая          & Средняя   \\ \hline
DAC    & Нет                        & Высокая  & Низкая           & Средняя   \\ \hline
MAC    & Нет                        & Низкая   & Средняя          & Высокая   \\ \hline
\end{tabular}
\end{table}

На основе проведенного анализа и сравнения различных моделей управления доступом можно выделить два наиболее подходящих подхода: RBAC и ABAC.

Пример таблицы с разрывом на две страницы \ref{tab:client}.

\begin{table}[H]
  \caption{Сравнение технологий для разработки клиентской части}
  \label{tab:client}
  \begingroup
    \begin{spacing}{1} % Одинарный интервал
      \begin{tabular}{|p{4.5cm}|p{3cm}|p{3cm}|p{4cm}|}
        \hline
        \multicolumn{1}{|c}{Характеристика}       & \multicolumn{1}{|c}{Next.js}                 & \multicolumn{1}{|c|}{Nuxt.js}                    & \multicolumn{1}{c|}{Angular}   \\ \hline
        \multicolumn{1}{|c}{1}       & \multicolumn{1}{|c}{2}                 & \multicolumn{1}{|c|}{3}                    & \multicolumn{1}{c|}{4} \\ \hline
        Базовая технология   & React                   & Vue.js                     & Фреймворк от Google         \\ \hline
        Язык разработки      & JavaScript / TypeScript & JavaScript / TypeScript    & TypeScript                  \\ \hline
      \end{tabular}
    \end{spacing}
  \endgroup
\end{table}

\begin{table}[H]
  \caption*{Продолжение таблицы \ref{tab:client}}
  \begingroup
    \begin{spacing}{1} % Одинарный интервал
      \begin{tabular}{|p{4.5cm}|p{3cm}|p{3cm}|p{4cm}|}
        \hline
        \multicolumn{1}{|c}{1}       & \multicolumn{1}{|c}{2}                 & \multicolumn{1}{|c|}{3}                    & \multicolumn{1}{c|}{4} \\ \hline
        Архитектура          & Компонентная            & Компонентная               & Модульная                   \\ \hline
        Сложность            & Умеренная               & Умеренная                  & Высокая                     \\ \hline
        Производительность   & Высокая                 & Высокая                    & Средняя на больших проектах \\ \hline
        Размер сообщества    & Большой                 & Большой, но меньше Next.js & Большое                     \\ \hline
        Скорость разработки  & Высокая                 & Высокая                    & Средняя                     \\ \hline
        Гибкость             & Высокая                 & Высокая                    & Ограниченная                \\ \hline
        Поддержка TypeScript & Отличная                & Хорошая                    & Нативная                    \\ \hline
      \end{tabular}
    \end{spacing}
  \endgroup
\end{table}
\vspace{10pt}

\section{Вывод}

Выводы по главе